\section{Nested gates and exact sequential elimination}
\label{sec:tpc62-nested-gates}

Transitivity for nested shorts is part of the classical shorted-operator
calculus \citep{Anderson1971,AndersonTrapp1975}.  The self-contained proof
below records precisely how that identity acts on the staged nuisance
spaces of the present closure problem.

The master formula eliminates the full parent-null space at once.  A
proof architecture may instead expose nuisance directions in several
layers.  This section shows that exact shorting is transitive along nested
retained spaces, so the sequential and one-shot procedures agree.  The
statement is purely finite-dimensional and makes no arithmetic assertion.

\subsection{Transitivity of the short}

\begin{theorem}[Transitivity on nested retained spaces]
\label{thm:tpc62-short-transitivity}
Let \(G\ge0\) on a finite-dimensional Hilbert space \(\cK\), and let
\begin{equation}
 T\subseteq S\subseteq\cK
 \label{eq:tpc62-nested-two-spaces}
\end{equation}
be subspaces.  Then
\begin{equation}
 \boxed{
 \Sh_T\!\bigl(\Sh_S(G)\bigr)=\Sh_T(G).}
 \label{eq:tpc62-short-transitivity}
\end{equation}
\end{theorem}

\begin{proof}
Use the maximal characterization
\eqref{eq:tpc62-short-order-definition}.  The operator
\(\Sh_T(\Sh_S(G))\) is supported on \(T\) and satisfies
\[
 0\le\Sh_T(\Sh_S(G))\le\Sh_S(G)\le G.
\]
It is therefore admissible in the defining maximum for \(\Sh_T(G)\), and
hence
\begin{equation}
 \Sh_T(\Sh_S(G))\le\Sh_T(G).
 \label{eq:tpc62-transitivity-first-order}
\end{equation}
Conversely, \(\Sh_T(G)\) is supported on \(T\subseteq S\) and is bounded
above by \(G\).  It is therefore admissible in the defining maximum for
\(\Sh_S(G)\), so \(\Sh_T(G)\le\Sh_S(G)\).  Since it is supported on
\(T\), it is then admissible in the defining maximum for
\(\Sh_T(\Sh_S(G))\).  Thus
\begin{equation}
 \Sh_T(G)\le\Sh_T(\Sh_S(G)).
 \label{eq:tpc62-transitivity-second-order}
\end{equation}
Combining the two inequalities proves
\eqref{eq:tpc62-short-transitivity}.
\end{proof}

There is also a direct variational reading.  Since \(T\subseteq S\),
\begin{equation}
 T^\perp=(S\cap T^\perp)\oplus S^\perp.
 \label{eq:tpc62-nested-complement-decomposition}
\end{equation}
For \(t\in T\), applying \eqref{eq:tpc62-short-variational} twice gives
\begin{align}
 \left\langle\Sh_T(\Sh_S(G))t,t\right\rangle
 &=\min_{u\in S\cap T^\perp}
   \min_{n\in S^\perp}
   \langle G(t+u+n),t+u+n\rangle
 \notag\\
 &=\min_{v\in T^\perp}
   \langle G(t+v),t+v\rangle
 \notag\\
 &=\langle\Sh_T(G)t,t\rangle.
 \label{eq:tpc62-transitivity-variational}
\end{align}
Thus transitivity is exactly the associativity of eliminating independent
orthogonal nuisance layers.  Nestedness is essential to this statement;
no commutation claim is made for two unrelated retained spaces.

\subsection{A ladder of nuisance spaces}

Let
\begin{equation}
 \cK=S_0\supseteq S_1\supseteq\cdots\supseteq S_k
 \label{eq:tpc62-retained-ladder}
\end{equation}
be a nested ladder, and define the nuisance layer removed at stage \(j\)
by
\begin{equation}
 E_j:=S_{j-1}\cap S_j^\perp,
 \qquad 1\le j\le k.
 \label{eq:tpc62-nuisance-layers}
\end{equation}
Then
\begin{equation}
 S_{j-1}=S_j\oplus E_j,
 \qquad
 S_j^\perp=E_j\oplus E_{j-1}\oplus\cdots\oplus E_1
 \qquad(1\le j\le k).
 \label{eq:tpc62-ladder-orthogonal-decomposition}
\end{equation}
Let \(F:\cK\to\cH\), set \(G_0=F^*F\), and recursively define
\begin{equation}
 G_j:=\Sh_{S_j}(G_{j-1}),
 \qquad 1\le j\le k.
 \label{eq:tpc62-sequential-shorts}
\end{equation}

\begin{theorem}[Sequential elimination equals one short]
\label{thm:tpc62-sequential-equals-direct}
For every \(j\),
\begin{equation}
 \boxed{G_j=\Sh_{S_j}(F^*F).}
 \label{eq:tpc62-sequential-direct-operator}
\end{equation}
Moreover, for \(s\in S_j\),
\begin{align}
 \langle G_js,s\rangle
 &=\min_{e_j\in E_j}\min_{e_{j-1}\in E_{j-1}}
   \cdots\min_{e_1\in E_1}
   \left\|F\!\left(s+\sum_{i=1}^j e_i\right)\right\|^2
 \label{eq:tpc62-sequential-nuisance-formula}\\
 &=\min_{n\in S_j^\perp}\|F(s+n)\|^2.
 \label{eq:tpc62-one-shot-nuisance-formula}
\end{align}
All minima are attained.
\end{theorem}

\begin{proof}
The operator identity follows inductively from
\cref{thm:tpc62-short-transitivity}:
\[
 G_j
 =\Sh_{S_j}\!\bigl(\Sh_{S_{j-1}}(F^*F)\bigr)
 =\Sh_{S_j}(F^*F).
\]
For the form identity, the variational formula for shorting from
\(S_{j-1}\) to \(S_j\) is
\begin{equation}
 \langle G_js,s\rangle
 =\min_{e_j\in E_j}
   \langle G_{j-1}(s+e_j),s+e_j\rangle.
 \label{eq:tpc62-one-stage-elimination}
\end{equation}
Iterating \eqref{eq:tpc62-one-stage-elimination} proves
\eqref{eq:tpc62-sequential-nuisance-formula}.  The orthogonal direct sum
\eqref{eq:tpc62-ladder-orthogonal-decomposition} identifies
\(\sum_{i=1}^j e_i\) with an arbitrary vector of \(S_j^\perp\), proving
\eqref{eq:tpc62-one-shot-nuisance-formula}.  Every space is finite
dimensional, and the least-squares argument
\eqref{eq:tpc62-short-least-squares} shows that the minima are attained.
\end{proof}

For a nonzero parent form in \cref{sec:tpc62-master-shorted}, take the
first retained space to be
\begin{equation}
 S_1=(\Ker R)^\perp.
 \label{eq:tpc62-parent-first-retained-space}
\end{equation}
The first eliminated layer is \(E_1=\Ker R\), and
\(G_1=\Sh_{S_1}(F^*F)\) is precisely the numerator in
\eqref{eq:tpc62-master-global-eigenvalue}.  If a subsequent argument has
genuinely nested retained spaces
\(S_k\subseteq\cdots\subseteq S_1\), then
\begin{equation}
 \Sh_{S_k}\!\left(\cdots
   \Sh_{S_2}\!\left(\Sh_{S_1}(F^*F)\right)\cdots\right)
 =\Sh_{S_k}(F^*F).
 \label{eq:tpc62-parent-all-at-once}
\end{equation}
Consequently no extra loss is created merely by staging exact nuisance
elimination.  Any loss appearing only after replacing a short by an upper
or lower certificate belongs to that certificate, not to the underlying
finite-dimensional elimination.

\subsection{Why compression is not elimination}

Ordinary compression to \(S\) is
\begin{equation}
 \cC_S(G):=P_SGP_S.
 \label{eq:tpc62-compression-definition}
\end{equation}
It evaluates the form after setting the discarded coefficient to zero.
Shorting instead minimizes over that coefficient.  Therefore
\begin{equation}
 0\le\Sh_S(G)\le\cC_S(G)
 \label{eq:tpc62-short-below-compression}
\end{equation}
as operators supported on \(S\), and equality need not hold.  The next
example shows that one compression in a nested chain can turn an exact
zero into a false positive constant.

\begin{example}[Intermediate compression destroys exact cancellation]
\label{ex:tpc62-nested-compression-failure}
Let \(\cK=\C^3\), \(\cH=\C^2\), and
\begin{equation}
 F(x,y,z)=(x+y,\,y+z).
 \label{eq:tpc62-compression-counterexample-map}
\end{equation}
Take
\begin{equation}
 S_1=\operatorname{span}\{e_1,e_2\},
 \qquad S_2=\operatorname{span}\{e_1\}.
 \label{eq:tpc62-compression-counterexample-spaces}
\end{equation}
The Gram matrix is
\begin{equation}
 G=F^*F=
 \begin{pmatrix}
 1&1&0\\
 1&2&1\\
 0&1&1
 \end{pmatrix}.
 \label{eq:tpc62-compression-counterexample-Gram}
\end{equation}
Exact elimination of \(z\) gives, on \(S_1\),
\begin{equation}
 \min_{z\in\C}\bigl(|x+y|^2+|y+z|^2\bigr)
 =|x+y|^2,
 \label{eq:tpc62-compression-counterexample-first-short}
\end{equation}
and exact elimination of \(y\) then gives zero.  Equivalently,
\begin{equation}
 \Sh_{S_2}(\Sh_{S_1}(G))
 =\Sh_{S_2}(G)=0,
 \label{eq:tpc62-compression-counterexample-exact-zero}
\end{equation}
because for any \(x\) one may choose \(y=-x\) and \(z=x\), making
\(F(x,y,z)=0\).

If the first short is incorrectly replaced by compression, the form on
\(S_1\) is instead
\begin{equation}
 |x+y|^2+|y|^2.
 \label{eq:tpc62-compression-counterexample-compressed-form}
\end{equation}
Shorting this compressed form to \(S_2\) yields
\begin{equation}
 \min_{y\in\C}\bigl(|x+y|^2+|y|^2\bigr)
 =\frac12|x|^2,
 \label{eq:tpc62-compression-counterexample-false-half}
\end{equation}
where the minimizer is \(y=-x/2\).  Thus
\begin{equation}
 \Sh_{S_2}(\cC_{S_1}(G))=\frac12 P_{S_2}\ne0.
 \label{eq:tpc62-compression-counterexample-operator-half}
\end{equation}
Compressing at both stages is still worse:
\begin{equation}
 \cC_{S_2}(\cC_{S_1}(G))=P_{S_2}.
 \label{eq:tpc62-compression-counterexample-operator-one}
\end{equation}
The numbers \(0\), \(1/2\), and \(1\) therefore arise from, respectively,
exact elimination, one premature compression, and two premature
compressions of the same positive Gram matrix.
\end{example}

The mechanism is not a numerical pathology.  In the example, the
discarded coordinate \(z\) cancels the second output component.  Setting
\(z=0\) retains a spurious penalty \(|y|^2\), which prevents the next
layer from using \(y=-x\) and manufactures a positive lower bound.  In a
physical closure argument, a coefficient direction may be compressed
away only after proving that it cannot alter the final output, or after
replacing the desired statement by one explicitly restricted to the
compressed coefficient class.

\subsection{Exact decisions and route failures}

Combining
\cref{thm:tpc62-master-shorted-gram,thm:tpc62-sequential-equals-direct}
gives a quantifier-safe decision rule.
After all genuine nuisance layers have been shorted, an actual zero
eigenvalue is equivalent to a final kernel vector carrying positive parent
energy.  It kills the uniform closure gate on the declared class.  By
contrast, any of the following is only a failure of a proposed proof
route:
\begin{enumerate}[label=\textup{(\roman*)},leftmargin=2.7em]
 \item replacing a shorted operator by a compression and obtaining an
       unusable bound;
 \item failing to control a sufficient singular-value or norm
       certificate;
 \item exhausting an exponent budget in an upper-bound calculation;
 \item proving a pointwise estimate for one distinguished vector but not
       the uniform generalized eigenvalue.
\end{enumerate}
None of these statements supplies a vector in
\(\Ker F\setminus\Ker R\).  Conversely, the mandatory-factor witness in
\cref{prop:tpc62-mandatory-factor-kill} does supply such a vector and
therefore survives every reformulation of the certificate.  This is the
strict L0 boundary used below: exact operator kernels decide impossibility
of the declared uniform gate, whereas failed certificates decide only
which route should be abandoned or redesigned.
